% ============================================================================
% Section: Data Sources and Data Generation
% ----------------------------------------------------------------------------
% Assumes: \usepackage[utf8]{inputenc}, \usepackage[T1]{fontenc}
%          \usepackage{booktabs}, \usepackage{array}
% (booktabs gives \toprule/\midrule/\bottomrule; array enables the p{} column
%  type used below for wrapped description columns; both are standard in
%  virtually all modern LaTeX distributions.)
% Adjust \label{...} keys if they collide with existing labels in the report.
% ============================================================================

\section{Data Sources and Data Generation}

This section describes the data foundation of the fraud-detection pipeline. The base transactional dataset is first characterized, including empirical properties discovered through preliminary profiling that directly shaped the subsequent design decisions. The methodology used to generate twelve additional behavioral and contextual fields is then presented, together with the validation procedure used to ensure these fields are realistic without leaking the fraud label. Finally, the accompanying data dictionary that documents every generated field is described.

\subsection{Base Dataset}
\label{subsec:base-dataset}

The base transactional data is drawn from the \emph{Online Payments Fraud Detection Dataset} (commonly referred to as PaySim), a publicly available Kaggle dataset generated by an agent-based mobile-money simulator. The simulator models a population of customers and merchants performing financial transactions over a one-month period, with a small subset of transactions labeled as fraudulent to emulate account-takeover behavior (an attacker gaining unauthorized access to an account and subsequently transferring or withdrawing funds).

The dataset comprises \textbf{6,362,620 transaction records} across \textbf{11 columns}, summarized in Table~\ref{tab:paysim-columns}. The full dataset was used without sampling throughout this work.

\begin{table}[htbp]
\centering
\caption{Base columns of the PaySim transactional dataset}
\label{tab:paysim-columns}
\footnotesize
\begin{tabular}{@{}llp{7.5cm}@{}}
\toprule
\textbf{Column} & \textbf{Type} & \textbf{Description} \\
\midrule
\texttt{step} & Integer & Elapsed simulation time, in hours, since the start of the simulation (range 1--743, i.e.\ approximately 31 days). \\
\texttt{type} & Categorical & Transaction type: \texttt{CASH\_IN}, \texttt{CASH\_OUT}, \texttt{DEBIT}, \texttt{PAYMENT}, or \texttt{TRANSFER}. \\
\texttt{amount} & Float & Transaction value, in the simulator's local currency unit. \\
\texttt{nameOrig} & String & Identifier of the account initiating the transaction (prefix \texttt{C}). \\
\texttt{oldbalanceOrg} & Float & Origin account balance immediately before the transaction. \\
\texttt{newbalanceOrig} & Float & Origin account balance immediately after the transaction. \\
\texttt{nameDest} & String & Identifier of the recipient account (prefix \texttt{C} for a customer, \texttt{M} for a merchant). \\
\texttt{oldbalanceDest} & Float & Destination account balance immediately before the transaction. \\
\texttt{newbalanceDest} & Float & Destination account balance immediately after the transaction. \\
\texttt{isFraud} & Binary & Ground-truth fraud label (1 = fraudulent, 0 = legitimate). \\
\texttt{isFlaggedFraud} & Binary & A naive rule-based flag built into the simulator (large-value transfers); present for only 16 of the 6,362,620 records and distinct from the true fraud label \texttt{isFraud}. \\
\bottomrule
\end{tabular}
\end{table}

The dataset exhibits the severe class imbalance typical of real-world fraud problems: only \textbf{8,213 of 6,362,620 transactions (0.1291\%)} are labeled fraudulent.

Before designing the synthetic-data-generation procedure described in Section~\ref{subsec:synthetic-data}, the base dataset was empirically profiled rather than assumed to behave in a particular way. Three findings from this profiling directly determined the subsequent design:

\begin{enumerate}
    \item \textbf{Fraud is confined to two transaction types.} Fraudulent transactions occur exclusively within \texttt{TRANSFER} (4,097 of 532,909 transactions) and \texttt{CASH\_OUT} (4,116 of 2,237,500 transactions); the \texttt{PAYMENT}, \texttt{CASH\_IN}, and \texttt{DEBIT} types contain zero fraudulent records. This indicates that the fraud modeled in this dataset represents \emph{account-takeover} behavior (unauthorized transfer or cash-out from a compromised account) rather than point-of-sale or card-not-present payment fraud.
    \item \textbf{Customer transaction history is negligible.} 6,353,307 of 6,362,620 distinct \texttt{nameOrig} identifiers (99.85\%) appear exactly once in the dataset. Consequently, features requiring a persistent multi-transaction customer profile would be constructed on essentially non-existent history; all synthetic behavioral fields were instead generated at the individual-transaction level.
    \item \textbf{The \texttt{step} field permits exact temporal derivation.} Because \texttt{step} encodes elapsed hours from a fixed simulation start with no gaps across its 1--743 range, time-of-day features can be computed deterministically rather than sampled stochastically.
\end{enumerate}

\subsection{Synthetic Contextual Data}
\label{subsec:synthetic-data}

The base dataset contains only raw financial transaction fields and lacks the behavioral and e-commerce contextual signals ordinarily available to a production fraud-detection system (device fingerprinting, geolocation consistency, account tenure, payment-retry behavior, and temporal usage patterns). To address this gap, twelve additional contextual fields were generated in Python using the \texttt{Faker} library in combination with custom, explicitly documented business logic.

\paragraph{Generation methodology.} All fields were generated through fully vectorized \texttt{numpy}/\texttt{pandas} operations applied to the complete dataset (no row-wise iteration, and no sub-sampling), drawing all randomness from a single seeded generator (\texttt{numpy.random.Generator}, seed \(=42\)) to guarantee exact reproducibility across runs. Five design principles governed the generation process:

\begin{enumerate}
    \item \textbf{Row-level generation.} Given the negligible customer transaction history established in Section~\ref{subsec:base-dataset}, fields were generated independently per transaction rather than through a persistent customer profile carried across multiple records.
    \item \textbf{Selective conditioning on the fraud label.} Only fields with a plausible behavioral connection to account-takeover fraud (e.g., use of an unrecognized device, a geographically inconsistent IP address, an unusually recent account, an atypical transaction hour, repeated failed payment attempts) were sampled conditionally on \texttt{isFraud}. Fields lacking such a rationale (e.g., browser, device type, device identifier, billing country) were generated independently of the label, so as not to manufacture an artificially perfect separation between fraudulent and legitimate transactions.
    \item \textbf{Bounded effect sizes.} Every conditional parameter (odds ratio, Poisson rate, or distributional median shift) was capped between two and four times its baseline value and documented as an explicit business assumption, rather than calibrated to any external ground truth (no such ground truth is publicly available for this dataset).
    \item \textbf{Deterministic derivation where possible.} Fields computable directly from existing columns (time-of-day features, and the geographic distance between IP and billing country) were derived by closed-form formula rather than sampled stochastically.
    \item \textbf{Post-hoc leakage validation.} After generation, every new field's univariate association with the fraud label was measured objectively (methodology described below), rather than relying on a subjective judgment of plausibility.
\end{enumerate}

Table~\ref{tab:synthetic-fields} summarizes the twelve generated fields, their generation mechanism, and a short description.

\begin{table}[htbp]
\centering
\caption{Synthetic contextual fields added to the base dataset}
\label{tab:synthetic-fields}
\footnotesize
\begin{tabular}{@{}llp{7cm}@{}}
\toprule
\textbf{Field} & \textbf{Mechanism} & \textbf{Description} \\
\midrule
\texttt{hour\_of\_day} & Derived & Hour of day (0--23), computed as \((\texttt{step}-1) \bmod 24\). \\
\texttt{is\_night\_transaction} & Derived & Boolean flag for transactions occurring between 00:00 and 05:59. \\
\texttt{customer\_account\_age\_days} & Conditional (log-normal) & Age of the customer account in days; median 400 days for legitimate transactions versus 275 days for fraudulent ones, reflecting the assumption that compromised/mule accounts skew more recently created. \\
\texttt{device\_id} & Independent & Device identifier sampled uniformly from a fixed pool of 50,000 Faker-generated UUIDs. \\
\texttt{browser} & Independent & Browser family (Chrome 55\%, Safari 20\%, Edge 12\%, Firefox 8\%, Other 5\%). \\
\texttt{device\_type} & Independent & Device category (mobile 65\%, desktop 30\%, tablet 5\%). \\
\texttt{new\_device\_flag} & Conditional (Bernoulli) & Whether the transaction originates from an unrecognized device; base rate 4\%, rising to 12\% for fraudulent transactions. \\
\texttt{billing\_country} & Independent & Customer's registered country, sampled from a fixed weighted distribution over 20 countries. \\
\texttt{ip\_country} & Conditional (Bernoulli) & Country inferred from the transaction's IP address; matches \texttt{billing\_country} with probability 93\% for legitimate transactions and 80\% for fraudulent ones. \\
\texttt{ip\_billing\_distance\_km} & Derived & Great-circle (haversine) distance between the \texttt{ip\_country} and \texttt{billing\_country} centroids; exactly zero when the two countries coincide. \\
\texttt{shipping\_billing\_mismatch} & Conditional (Bernoulli) & Whether the transaction's associated address differs from the registered address; base rate 5\%, rising to 15\% for fraudulent transactions. \\
\texttt{failed\_payment\_attempts\_24h} & Conditional (Poisson) & Number of failed payment attempts in the preceding 24 hours; mean 0.15 for legitimate transactions, rising to 0.6 for fraudulent ones. \\
\bottomrule
\end{tabular}
\end{table}

\paragraph{Leakage validation.} A synthetic feature that correlates too strongly with the fraud label would trivialize the downstream classification task and misrepresent the value of the generated context. Each field was therefore validated after generation: for continuous or binary fields, the univariate area under the ROC curve (AUC) against \texttt{isFraud} was computed; for categorical fields, a bias-corrected Cram\'er's V statistic (Bergsma and Wicher's small-sample correction) was used in place of the naive formula, since the naive estimator is known to be upward-biased on sparse contingency tables and was found, during independent verification, to be sensitive to the cardinality of the \texttt{device\_id} field and to dataset size. A field was flagged as failing the leakage check if its AUC reached or exceeded 0.75, or its corrected Cram\'er's V reached or exceeded 0.50.

The validation procedure was applied iteratively rather than treated as a single pass. On the first run against the full 6,362,620-row dataset, \texttt{customer\_account\_age\_days} exceeded the AUC threshold (0.8753); the fraud-conditional median was consequently reduced from 150 to 275 days (narrowing, rather than eliminating, the intended signal) and the field was regenerated, after which its AUC fell to 0.6689. A subsequent independent review additionally identified that the originally implemented (uncorrected) Cram\'er's V formula could report a spuriously high, sample-size-dependent value for \texttt{device\_id} despite that field being generated independently of the fraud label by construction; substituting the bias-corrected estimator resolved this without any further change to the generation logic. Table~\ref{tab:leakage-results} reports the final validation results on the complete dataset; all twelve fields satisfy the leakage criteria.

\begin{table}[htbp]
\centering
\caption{Final leakage-validation results on the full dataset (6,362,620 rows)}
\label{tab:leakage-results}
\footnotesize
\begin{tabular}{@{}lccc@{}}
\toprule
\textbf{Field} & \textbf{Metric} & \textbf{Value} & \textbf{Status} \\
\midrule
\texttt{hour\_of\_day} & AUC & 0.6336 & PASS \\
\texttt{is\_night\_transaction} & AUC & 0.6217 & PASS \\
\texttt{customer\_account\_age\_days} & AUC & 0.6689 & PASS \\
\texttt{device\_id} & Corrected Cram\'er's V & 0.0000 & PASS \\
\texttt{browser} & Corrected Cram\'er's V & 0.0000 & PASS \\
\texttt{device\_type} & Corrected Cram\'er's V & 0.0004 & PASS \\
\texttt{new\_device\_flag} & AUC & 0.5419 & PASS \\
\texttt{billing\_country} & Corrected Cram\'er's V & 0.0000 & PASS \\
\texttt{ip\_country} & Corrected Cram\'er's V & 0.0049 & PASS \\
\texttt{ip\_billing\_distance\_km} & AUC & 0.5651 & PASS \\
\texttt{shipping\_billing\_mismatch} & AUC & 0.5495 & PASS \\
\texttt{failed\_payment\_attempts\_24h} & AUC & 0.6589 & PASS \\
\bottomrule
\end{tabular}
\end{table}

Generating the twelve fields did not alter the underlying class distribution: the resulting dataset retains all 6,362,620 records and the original 0.1291\% fraud rate, since rebalancing (e.g.\ oversampling or class weighting) is deliberately deferred to the feature-engineering and modeling stage.

\subsection{Data Dictionary}
\label{subsec:data-dictionary}

In accordance with the requirement that every synthetically generated field be formally documented, a structured data dictionary was produced covering all twelve fields listed in Table~\ref{tab:synthetic-fields}. For each field, the dictionary records: the column name, data type, unit, valid range, generation mechanism (derived, independently sampled, or conditional on the fraud label), the exact generation formula or distribution, the underlying business assumption, and the measured univariate association with \texttt{isFraud} reported in Table~\ref{tab:leakage-results}.

Rather than being transcribed by hand, the data dictionary is generated programmatically by the same validation module that performs the leakage checks described in Section~\ref{subsec:synthetic-data}. This guarantees that the documented specification and the executed generation code cannot drift out of sync: any change to a generation parameter is automatically reflected in the published dictionary the next time the pipeline is run, and every reported association value is a direct measurement on the current dataset rather than a manually maintained estimate. The dictionary is published in Markdown format, satisfying the required deliverable format (Markdown or Excel). Table~\ref{tab:data-dict-excerpt} illustrates its structure with three representative entries.

\begin{table}[htbp]
\centering
\caption{Excerpt of the generated data dictionary (illustrative entries)}
\label{tab:data-dict-excerpt}
\footnotesize
\begin{tabular}{@{}lp{2.2cm}p{2.2cm}p{6cm}@{}}
\toprule
\textbf{Field} & \textbf{Type / Range} & \textbf{Generation} & \textbf{Business assumption} \\
\midrule
\texttt{new\_device\_flag} & Boolean & Conditional (Bernoulli, $p=0.04\!\to\!0.12$) & Account-takeover fraud is plausibly more often initiated from an unrecognized device, though not certainly so, since sophisticated attackers may spoof a known device fingerprint. \\
\texttt{ip\_billing\_distance\_km} & Float, [0, $\sim$20{,}000] km & Derived from \texttt{ip\_country} / \texttt{billing\_country} & Computed directly from a fixed country-centroid lookup table to guarantee internal consistency with the country-mismatch fields. \\
\texttt{failed\_payment\_attempts\_24h} & Integer, [0, $\sim$5] & Conditional (Poisson, $\lambda=0.15\!\to\!0.6$) & Attackers are assumed to make repeated failed attempts (e.g.\ credential or card guessing) before a successful transaction. \\
\bottomrule
\end{tabular}
\end{table}
